\chapter{Full Working Examples}
\label{app:worked}

This appendix collects complete, working source listings for the games built and
referenced in the chapters. Each is reproduced verbatim from the \prg repository's
\fname{examples/games} directory, where the sources are deliberately written with a
comment on nearly every instruction so that a first-year student can connect each
line to its visible effect \citep{prg32repo}. Read them after the corresponding
chapter: the assembly listings reward the reader of Part~II, the C listings the
reader of Part~III, and the side-by-side comparison rewards anyone who has read
Chapter~\ref{ch:join}.

All four programs implement games discussed in the text, and all four are designed
to run unchanged on both the QEMU emulator and the physical ESP32\nobreakdash-C6, packaged as
\fname{.prg32} cartridges by the workflow of Appendix~\ref{app:api}
\citep{prg32qemu}.

\section{Pong in graphics assembly}

This is the complete \code{pong\_graphics} example developed across
Chapters~\ref{ch:asm1}--\ref{ch:asm3}. It exports the three symbols
\code{pong\_graphics\_init}, \code{pong\_graphics\_update}, and
\code{pong\_graphics\_draw}, keeps its state in \code{.data}, moves a paddle and a
bouncing ball, and renders them with \code{prg32\_gfx\_rect}. The optional debug
overlay at the end is guarded by a compile-time flag.

\lstinputlisting[style=asm,caption={\fname{examples/games/pong/graphics/game.S} --- complete graphics Pong in \riscv assembly \citep{prg32repo}.}]{examples/pong_graphics_game.S}

\section{Pong in ASCII assembly}

The ASCII variant uses the same three-function structure but renders to the text
console, making it the gentlest possible introduction to the calling convention,
as used in Chapter~\ref{ch:asm1}. Note how closely its \code{update} mirrors the
graphics version: the same input read, the same bit tests, the same clamps,
differing only in units (text columns rather than pixels).

\lstinputlisting[style=asm,caption={\fname{examples/games/pong/ascii/game.S} --- complete ASCII Pong in \riscv assembly \citep{prg32repo}.}]{examples/pong_ascii_game.S}

\section{Pong in C}

The C version, the basis of Chapter~\ref{ch:c1}, expresses the identical game in a
fraction of the lines. Compare it directly with the graphics assembly above: every
C statement here corresponds to a block of instructions there, which is the lesson
of Chapter~\ref{ch:join}. This version also keeps a second paddle alive with a
fallback AI so the two-player screen stays in motion even without a second player.

\lstinputlisting[style=c,caption={\fname{examples/games/pong/c/game.c} --- complete Pong in C \citep{prg32repo}.}]{examples/pong_c_game.c}

\section{Breakout in C}

Breakout is Pong with one new idea: a row of bricks, stored as the bits of a single
byte, each set bit meaning a brick is still present. When the ball reaches the
brick row, the code clears the corresponding bit with a bitwise AND and bounces the
ball. This is the bitmask technique of Chapter~\ref{ch:asm2} reused for game state,
as discussed in Chapter~\ref{ch:asmcap}.

\lstinputlisting[style=c,caption={\fname{examples/games/breakout/c/game.c} --- complete Breakout in C, showing bricks as a bitmask \citep{prg32repo}.}]{examples/breakout_c_game.c}

\section{Where to find the rest}

The repository contains, for each of the nine example games, an ASCII assembly
version, a graphics assembly version, and a C version, alongside focused feature
demos for scrolling, parallax, animated sprites, and dual playfields
\citep{prg32repo}. Having read the four listings above, you have the vocabulary to
read any of them. The recommended next reads, in increasing richness, are Space
Invaders, Tetris, the platformer, the raycaster, and Wing Commander, each of which
is a variation on the patterns reproduced here.
